\documentclass[11pt, a4paper]{article}

\usepackage[utf8]{inputenc}
\usepackage[T1]{fontenc}
\usepackage[ngerman]{babel}
\usepackage{amsmath, amssymb}
\usepackage{booktabs}
\usepackage{listings}
\usepackage{xcolor}
\usepackage{hyperref}
\usepackage{geometry}
\usepackage{parskip}
\usepackage{microtype}

\geometry{margin=2.5cm}

\definecolor{codebg}{gray}{0.95}
\definecolor{codecomment}{rgb}{0.4,0.4,0.4}
\definecolor{codestring}{rgb}{0.13,0.54,0.13}
\definecolor{codekeyword}{rgb}{0.0,0.0,0.6}

\lstset{
  language=Python,
  backgroundcolor=\color{codebg},
  basicstyle=\ttfamily\small,
  keywordstyle=\color{codekeyword}\bfseries,
  commentstyle=\color{codecomment}\itshape,
  stringstyle=\color{codestring},
  breaklines=true,
  frame=single,
  framerule=0pt,
  rulecolor=\color{codebg},
  xleftmargin=1em,
  xrightmargin=1em,
  aboveskip=0.8em,
  belowskip=0.8em,
}

\lstdefinestyle{sql}{
  language=SQL,
  backgroundcolor=\color{codebg},
  basicstyle=\ttfamily\small,
  keywordstyle=\color{codekeyword}\bfseries,
  breaklines=true,
  frame=single,
  framerule=0pt,
  rulecolor=\color{codebg},
  xleftmargin=1em,
  xrightmargin=1em,
}

\title{\textbf{Pipeline-Dokumentation: USD Risk-Free Rate Curve}\\
       \large Projekt \texttt{research\_db} --- Modul 5}
\author{Research-Datenbank-Projekt}
\date{\today}

% ============================================================
\begin{document}
\maketitle
\tableofcontents
\newpage

% ============================================================
\section{Zweck und Einordnung}

Dieses Dokument beschreibt die Implementierung und das Design der
USD-Zinskurven-Pipeline im Rahmen des \texttt{research\_db}-Projekts.

Die tägliche USD Risk-Free Rate Curve liefert die Grundlage für die
Diskontierung von Cashflows in der stocklevel SVIX-Pipeline
\cite{martinwagner2019}. Konkret wird die Kurve für folgende Zwecke
verwendet:

\begin{itemize}
  \item Diskontierung im CRR/Leisen-Reimer Binomialbaum zur Berechnung
        europäisch-äquivalenter Optionspreise.
  \item Interpolation des risikofreien Zinssatzes für beliebige
        Optionslaufzeiten~$\tau$.
  \item Ersatz des bisherigen Hilfsverfahrens (SPX Put-Call-Parität als
        Proxy für den Diskontfaktor).
\end{itemize}

\subsection*{Projektpfade}

\begin{tabular}{ll}
  \toprule
  Komponente & Pfad \\
  \midrule
  Data-Projekt    & \texttt{C:\textbackslash Users\textbackslash Miebs\textbackslash PycharmProjects\textbackslash Data} \\
  Research-Projekt & \texttt{C:\textbackslash Users\textbackslash Miebs\textbackslash PycharmProjects\textbackslash Research} \\
  Datenbank       & PostgreSQL \texttt{research\_db} \\
  Modul           & \texttt{modules/fetch\_usd\_rates.py} \\
  \bottomrule
\end{tabular}

% ============================================================
\section{Datenbankschema}

\subsection{Tabellendefinition}

\begin{lstlisting}[style=sql]
CREATE SCHEMA IF NOT EXISTS pub_rates;

CREATE TABLE IF NOT EXISTS pub_rates.usd_rates (
    trade_date   DATE         NOT NULL,
    tenor        VARCHAR(8)   NOT NULL,
    tenor_days   SMALLINT     NOT NULL,
    rate_type    VARCHAR(8)   NOT NULL,
    rate_pct     NUMERIC(8,4),
    compounding  VARCHAR(16)  NOT NULL,
    source       VARCHAR(32)  NOT NULL,
    updated_at   TIMESTAMPTZ  NOT NULL DEFAULT now(),
    PRIMARY KEY (trade_date, tenor, rate_type)
);

CREATE INDEX IF NOT EXISTS idx_usd_rates_date
    ON pub_rates.usd_rates (trade_date);
\end{lstlisting}

\subsection{Feldbeschreibung}

\begin{tabular}{lll}
  \toprule
  Feld & Typ & Beschreibung \\
  \midrule
  \texttt{trade\_date}  & \texttt{DATE}        & Referenzdatum der Zinsbeobachtung \\
  \texttt{tenor}        & \texttt{VARCHAR(8)}  & Lesbares Tenor-Label (z.\,B.\ \texttt{'3m'}, \texttt{'10y'}) \\
  \texttt{tenor\_days}  & \texttt{SMALLINT}    & Tenor in Kalendertagen \\
  \texttt{rate\_type}   & \texttt{VARCHAR(8)}  & \texttt{'par'} oder \texttt{'zero'} \\
  \texttt{rate\_pct}    & \texttt{NUMERIC(8,4)}& Zinssatz in Prozent (z.\,B.\ $4{,}25$ für 4,25\,\%) \\
  \texttt{compounding}  & \texttt{VARCHAR(16)} & Verzinsungskonvention (siehe \S\ref{sec:conventions}) \\
  \texttt{source}       & \texttt{VARCHAR(32)} & \texttt{'FRED'} oder \texttt{'bootstrapped'} \\
  \texttt{updated\_at}  & \texttt{TIMESTAMPTZ} & Zeitstempel des letzten Schreibvorgangs \\
  \bottomrule
\end{tabular}

\subsection{Tenors und FRED-Serien}

\begin{tabular}{lrrl}
  \toprule
  Tenor & Tage & FRED-Serie & Beschreibung \\
  \midrule
  \texttt{1m}  &    30 & \texttt{DGS1MO} & 1-Month Treasury CMT \\
  \texttt{3m}  &    91 & \texttt{DGS3MO} & 3-Month Treasury CMT \\
  \texttt{6m}  &   182 & \texttt{DGS6MO} & 6-Month Treasury CMT \\
  \texttt{1y}  &   365 & \texttt{DGS1}   & 1-Year Treasury CMT \\
  \texttt{2y}  &   730 & \texttt{DGS2}   & 2-Year Treasury CMT \\
  \texttt{3y}  &  1095 & \texttt{DGS3}   & 3-Year Treasury CMT \\
  \texttt{5y}  &  1825 & \texttt{DGS5}   & 5-Year Treasury CMT \\
  \texttt{7y}  &  2555 & \texttt{DGS7}   & 7-Year Treasury CMT \\
  \texttt{10y} &  3650 & \texttt{DGS10}  & 10-Year Treasury CMT \\
  \texttt{20y} &  7300 & \texttt{DGS20}  & 20-Year Treasury CMT \\
  \texttt{30y} & 10950 & \texttt{DGS30}  & 30-Year Treasury CMT \\
  \bottomrule
\end{tabular}

% ============================================================
\section{Zinssatzkonventionen}
\label{sec:conventions}

\subsection{Par-Yields (\texttt{rate\_type = 'par'})}

Die FRED Constant Maturity Treasury (CMT) Rates werden unverändert
gespeichert:

\begin{itemize}
  \item \textbf{Kurzläufer 1m, 3m, 6m}: Bank-Discount-Basis,
        annualisiert auf Basis eines 360-Tage-Jahres.
        \texttt{compounding = 'discount'}
  \item \textbf{Langläufer 1y--30y}: Semiannuale Bond Equivalent Yield
        (BEY). \texttt{compounding = 'semiannual'}
\end{itemize}

\subsection{Zero-Coupon-Rates (\texttt{rate\_type = 'zero'})}

Werden per Bootstrap aus den Par-Yields abgeleitet.
\texttt{compounding = 'annual'} für alle Zero-Rate-Zeilen.

Diskontfaktor:
\begin{equation}
  \mathrm{disc}(\tau) = \left(1 + \frac{r_{\mathrm{zero}}}{100}\right)^{-\tau/365}
  \label{eq:discfactor}
\end{equation}

Umrechnung in Application Code (diskret $\to$ kontinuierlich):
\begin{equation}
  r_{\mathrm{cont}} = \ln\!\left(1 + \frac{r_{\mathrm{disc}}}{100}\right)
  \label{eq:disc2cont}
\end{equation}

% ============================================================
\section{Bootstrap-Algorithmus}
\label{sec:bootstrap}

\subsection{Short-End: Bank-Discount zu Zero-Rate ($\leq 6$M)}

T-Bill-Sätze liegen als Bank-Discount-Rate $d$ (in Prozent) vor.
Umrechnung in jährlich diskrete Zero-Rate:

\begin{align}
  P &= 1 - \frac{d}{100} \cdot \frac{n}{360} \label{eq:tbill_price} \\[4pt]
  r_{\mathrm{zero}} &= \left(\frac{1}{P}\right)^{365/n} - 1 \label{eq:tbill_zero}
\end{align}

wobei $n$ die Laufzeit in Kalendertagen ist.

\subsection{Long-End: Sequentieller Bootstrap ($\geq 1$Y)}

Treasury Bonds mit semiannualer Kuponzahlung und Nennwert 100.
Par-Bonds handeln zum Kurs 100 (Kupon = Rendite).

Die Par-Bond-Bedingung lautet:
\begin{equation}
  100 = \frac{c}{2} \cdot 100 \cdot \sum_{i=1}^{n-1} d(t_i)
        + \left(1 + \frac{c}{2}\right) \cdot 100 \cdot d(T)
  \label{eq:parbond}
\end{equation}

wobei:
\begin{itemize}
  \item $c$ die jährliche Kuponrate (als Dezimalzahl)
  \item $t_i = i \cdot 0{,}5 \cdot 365$ Tage, $i = 1, \ldots, n{-}1$
  \item $n = \mathrm{round}(T_{\mathrm{years}} \cdot 2)$ die Anzahl der semiannualen Perioden
  \item $d(t_i)$ der Diskontfaktor aus bereits gebootstrappten Zero-Rates
\end{itemize}

Auflösen nach $d(T)$:
\begin{equation}
  d(T) = \frac{1 - \dfrac{c}{2} \displaystyle\sum_{i=1}^{n-1} d(t_i)}{1 + \dfrac{c}{2}}
  \label{eq:solve_dT}
\end{equation}

Zero-Rate aus Diskontfaktor:
\begin{equation}
  r_{\mathrm{zero}}(T) = d(T)^{-365/T_{\mathrm{days}}} - 1
  \label{eq:zero_from_d}
\end{equation}

\subsection{Interpolation von Zwischen-Cashflows}

Zwischentermine $t_i$, die nicht exakt auf einen gebootstrappten Tenor
fallen, werden log-linear interpoliert. Für $t_1 \leq t \leq t_2$:

\begin{equation}
  r(t) = (1-w)\,r(t_1) + w\,r(t_2),
  \qquad
  w = \frac{\ln t - \ln t_1}{\ln t_2 - \ln t_1}
  \label{eq:loglin}
\end{equation}

An den Rändern (außerhalb des bekannten Bereichs) wird flach
extrapoliert.

\subsection{Bootstrap-Reihenfolge}

Die Reihenfolge der Bootstrap-Schritte ist:

\begin{enumerate}
  \item \textbf{1m} --- Bank-Discount (Gl.~\ref{eq:tbill_price}--\ref{eq:tbill_zero})
  \item \textbf{3m} --- Bank-Discount
  \item \textbf{6m} --- Bank-Discount
  \item \textbf{1y} --- Par-Bond Bootstrap; Zwischenterm 0,5y = exakter 6m-Zero
  \item \textbf{2y} --- Par-Bond Bootstrap; 0,5y, 1,0y exakt; 1,5y interpoliert
  \item \textbf{3y} --- Par-Bond Bootstrap; Zwischenterme aus 1y/2y-Zeros
  \item \textbf{5y, 7y, 10y, 20y, 30y} --- analog
\end{enumerate}

% ============================================================
\section{Implementierung}

\subsection{Modulstruktur}

\begin{tabular}{ll}
  \toprule
  Funktion & Aufgabe \\
  \midrule
  \texttt{\_ensure\_schema\_and\_table()} & Idempotente DDL-Ausführung beim Start \\
  \texttt{\_get\_fetch\_start()}           & Inkrementelles Startdatum (MAX + 1 Tag) \\
  \texttt{\_interp\_zero()}               & Log-lineare Interpolation der Zero-Kurve \\
  \texttt{\_bill\_to\_zero()}             & Bank-Discount $\to$ annual discrete zero \\
  \texttt{\_bond\_to\_zero()}             & Par-Bond-Bootstrap für einen Tenor \\
  \texttt{\_bootstrap\_zeros()}           & Bootstrap für einen Handelstag (alle Tenors) \\
  \texttt{\_build\_records()}             & Datenbankzeilen für Par + Zero \\
  \texttt{run(snapshot\_date)}            & Haupteinstiegspunkt \\
  \bottomrule
\end{tabular}

\subsection{Inkrementelles Laden}

\begin{lstlisting}
def _get_fetch_start() -> date:
    with engine.connect() as conn:
        row = conn.execute(
            text("SELECT MAX(trade_date) FROM pub_rates.usd_rates")
        ).scalar()
    if row is not None:
        return row + timedelta(days=1)
    return date.today().replace(year=date.today().year - 10)
\end{lstlisting}

Bei leerem Table: 10 Jahre Backfill ab heutigem Datum.
Bei bestehendem Bestand: Start ab dem Tag nach dem letzten bekannten Handelstag.
Wochenenden und US-Feiertage werden automatisch übersprungen
(FRED liefert für diese Tage keine Beobachtungen).

\subsection{Idempotenz}

Schreibvorgänge nutzen \texttt{ON CONFLICT DO NOTHING} auf dem
Primärschlüssel \texttt{(trade\_date, tenor, rate\_type)}. Ein
wiederholter Lauf für denselben Zeitraum überschreibt keine Daten.

% ============================================================
\section{Verwendung in der SVIX-Pipeline}

\subsection{Abfrage und Interpolation}

Für einen Handelstag und eine Laufzeit $\tau$ (in Tagen) wird der
risikofreie Zinssatz wie folgt ermittelt:

\begin{lstlisting}
def interpolate_zero_rate(trade_date, tau, conn):
    rows = conn.execute(text("""
        SELECT tenor_days, rate_pct
        FROM pub_rates.usd_rates
        WHERE trade_date = :d AND rate_type = 'zero'
        ORDER BY tenor_days
    """), {"d": trade_date}).fetchall()

    days = [r[0] for r in rows]
    rates = [r[1] for r in rows]

    # log-linear interpolation (flat extrapolation at edges)
    if tau <= days[0]:
        return float(rates[0])
    if tau >= days[-1]:
        return float(rates[-1])
    for i in range(len(days) - 1):
        if days[i] <= tau <= days[i+1]:
            w = (np.log(tau) - np.log(days[i])) / \
                (np.log(days[i+1]) - np.log(days[i]))
            return (1-w)*float(rates[i]) + w*float(rates[i+1])
\end{lstlisting}

\subsection{Diskontfaktor für den Binomialbaum}

\begin{lstlisting}
# 1. Zero-Rate interpolieren (in Prozent, annual discrete)
r_zero_pct = interpolate_zero_rate(trade_date, tau, conn)

# 2. Diskret -> kontinuierlich
r_cont = np.log(1 + r_zero_pct / 100)

# 3. Diskontfaktor
disc = np.exp(-r_cont * tau / 365)
\end{lstlisting}

% ============================================================
\section{Betrieb und Automatisierung}

\subsection{Manueller Aufruf}

\begin{lstlisting}[language=bash]
cd C:\Users\Miebs\PycharmProjects\Data
python -m modules.fetch_usd_rates
\end{lstlisting}

\subsection{Integration in run\_monthly.py}

Das Modul ist als Schritt~5 in den monatlichen Pipeline-Orchestrator
eingebunden:

\begin{lstlisting}
# Modul 5: USD Risk-Free Rates
from modules.fetch_usd_rates import run as run_usd_rates
run_usd_rates(snapshot_date)
\end{lstlisting}

\subsection{Windows Task Scheduler}

Für die tägliche Aktualisierung ist ein separater Windows-Task
eingerichtet:

\begin{tabular}{ll}
  \toprule
  Parameter & Wert \\
  \midrule
  Task-Name          & \texttt{fetch\_usd\_rates} \\
  Zeitplan           & täglich, 22:00 Uhr \\
  Batch-Script       & \texttt{run\_usd\_rates.bat} \\
  Log-Datei          & \texttt{logs/usd\_rates.log} \\
  \texttt{StartWhenAvailable} & Ja (Nachholung bei abgeschaltetem PC) \\
  \bottomrule
\end{tabular}

Die Uhrzeit 22:00 (MEZ/MESZ) stellt sicher, dass FRED die
CMT-Rates bereits veröffentlicht hat (typisch: $\sim$21:30 MEZ,
entspricht $\sim$15:30 ET).

\subsection{Historischer Bestand (Stand: März 2026)}

\begin{tabular}{lr}
  \toprule
  Kennzahl & Wert \\
  \midrule
  Erster Handelstag  & 22.03.2016 \\
  Letzter Handelstag & 19.03.2026 \\
  Anzahl Handelstage & 2.498 \\
  Zeilen gesamt      & 54.956 \\
  Zeilen pro Tag     & 22 (11 Par + 11 Zero) \\
  \bottomrule
\end{tabular}

% ============================================================
\begin{thebibliography}{9}
\bibitem{martinwagner2019}
  Martin, I. \& Wagner, C. (2019).
  \textit{What is the Expected Return on a Stock?}
  The Journal of Finance, 74(4), 1887--1929.
\end{thebibliography}

\end{document}
